\chapter{Resultados}

Este capítulo separa el estado de la evidencia de su interpretación. En la fecha de este borrador, la campaña final definida en \path{experiments/final_experiment_campaign.json} todavía no se ha ejecutado. En consecuencia, primero se fija la estructura en la que se incorporarán sus artefactos y, después, se conserva una síntesis de resultados históricos útiles para supervisión. Estos últimos no constituyen la evaluación final del sistema y serán sustituidos o relegados cuando exista una campaña completa y validada.

\section{Condiciones generales y trazabilidad}
\label{sec:resultados-condiciones}

Una cifra se considera evidencia final únicamente si puede trazarse desde una ejecución de la campaña aprobada hasta su \texttt{config.json}, \texttt{metrics.json} o fichero de validación, \texttt{environment.json}, resumen, manifiesto y sumas de comprobación. La Figura~\ref{fig:trazabilidad-evidencia} resume esa cadena. Esta regla evita mezclar tres objetos distintos: el perfil experimental congelado, los artefactos históricos ya existentes y los artefactos finales aún pendientes.

\begin{figure}[htbp]
\centering
\resizebox{0.98\textwidth}{!}{% Figura conceptual: cadena de evidencia, sin resultados medidos.
% Fuente: esquema de artefactos 3.0 y agregador de campaña.
\begin{tikzpicture}[node distance=3mm]
  \node[cajaDato, text width=2.1cm] (spec) {Diseño congelado\\+ perfil hash};
  \node[cajaProceso, text width=2.1cm, right=of spec] (run) {Ejecución física\\+ RUN\_ID};
  \node[cajaProceso, text width=2.1cm, right=of run] (art) {Artefactos 3.0\\+ SHA-256};
  \node[cajaModelo, text width=2.1cm, right=of art] (agg) {Agregado\\validado};
  \node[cajaResaltada, text width=2.1cm, right=of agg] (thesis) {Tabla / figura /\\afirmación};

  \draw[flechaAzul] (spec) -- (run);
  \draw[flechaAzul] (run) -- (art);
  \draw[flechaAzul] (art) -- (agg);
  \draw[flechaAzul] (agg) -- (thesis);

  \node[cajaPendiente, text width=8.6cm, below=7mm of art] (gate)
    {\textbf{Puerta de aceptación:} campaña completa, hashes coherentes, dependencias correctas, métricas definidas o nulas y ausencia de sustitutos históricos};
  \draw[flechaFina] (art) -- (gate);
  \draw[flechaFina] (gate.east) -| (agg.south);
\end{tikzpicture}
}
\caption[Trazabilidad desde el diseño hasta la tesis]{Cadena de trazabilidad exigida para incorporar una cifra final: especificación congelada, ejecución identificada, artefactos sellados, agregación validada y elemento de tesis. Una ausencia o inconsistencia detiene la incorporación. Fuente: elaboración propia a partir del esquema de artefactos de campaña.}
\label{fig:trazabilidad-evidencia}
\end{figure}

La Tabla~\ref{tab:estructura-resultados-finales} fija el orden definitivo. No se crean tablas vacías para simular resultados: cada bloque se incorporará cuando el agregador haya comprobado que están presentes todas las ejecuciones físicas y alias previstos.

\begin{table}[htbp]
\centering
\footnotesize
\setlength{\tabcolsep}{4pt}
\begin{tabular}{p{3.0cm}p{5.1cm}p{4.4cm}}
\toprule
Bloque final & Pregunta que resuelve & Evidencia requerida \\
\midrule
MAIN y controles directos & ¿Aprende el perfil congelado en la partición aleatoria y coinciden evaluación y validación independiente? & MAIN nueva, validación directa, \emph{bootstrap}, duplicados y etiquetas barajadas \\
Partición completa por días & ¿Persiste el comportamiento al romper la mezcla aleatoria entre días? & QRDQN y RF con lunes--miércoles para entrenamiento y jueves--viernes para prueba \\
Escalera de tamaños & ¿Cómo cambia el resultado con la cantidad de datos y el presupuesto proporcional de pasos? & Seis puntos lógicos entre $100\,000$ filas y el conjunto completo \\
Sensibilidad a semilla & ¿Cuánto varía el modelo al mantener fija la partición? & Cinco semillas de modelo, 42--46, sobre $1\,000\,000$ de filas \\
Escenarios retenidos & ¿Qué ataques o capturas inducen una degradación específica? & Cuatro \emph{holdouts} dirigidos y sus RF equivalentes \\
Comparación de modelos & ¿Qué parte del comportamiento depende de QRDQN y cuál del protocolo tabular? & Comparaciones emparejadas QRDQN--RF \\
Fase 2 & ¿Puede el MAIN final procesar la captura de laboratorio y con qué cambio de dominio? & Inferencia fresca, métricas y diagnósticos de distribución \\
\bottomrule
\end{tabular}
\caption[Estructura de los resultados finales]{Orden de incorporación de la campaña final y contrato mínimo de evidencia para cada bloque. La tabla describe decisiones analíticas, no resultados ejecutados. Fuente: elaboración propia a partir de \path{experiments/final_experiment_campaign.json}.}
\label{tab:estructura-resultados-finales}
\end{table}

\evidencenote{Los siete bloques de la Tabla~\ref{tab:estructura-resultados-finales} permanecen sin valores finales. La actualización posterior deberá producirse desde el agregado validado de campaña, nunca mediante transcripción estimada o reutilización silenciosa de cifras históricas.}

\section{Evidencia histórica previa a la campaña final}
\label{sec:evidencia-historica}

Los siguientes resultados proceden de artefactos comprometidos antes de congelar la campaña actual. Siguen siendo útiles para comprobar el funcionamiento del pipeline, identificar amenazas y justificar el diseño de la campaña; no deben leerse como una ejecución anticipada de esta.

\subsection{MAIN histórica: capacidad de aprendizaje intradistribución}
\label{sec:resultados-main}

La ejecución histórica \path{MAIN_qrdqn_cicids2017_canonical_full_random_20260609_193655} entrenó QRDQN durante $3\,000\,000$ de pasos sobre una partición aleatoria estratificada con semilla~42. Sobre $566\,149$ filas de prueba registró una exactitud de $0.99381$, \emph{recall} de ataque de $0.99536$ y F1 de ataque de $0.98445$. La matriz de confusión fue $\mathrm{TN}=451\,631$, $\mathrm{FP}=2\,989$, $\mathrm{FN}=518$ y $\mathrm{TP}=111\,011$.

\begin{table}[htbp]
\centering
\begin{tabular}{lc}
\toprule
Métrica & MAIN histórica \\
\midrule
Exactitud & $0.99381$ \\
Precisión de ataque & $0.97378$ \\
Recall de ataque & $0.99536$ \\
F1 de ataque & $0.98445$ \\
Tasa de falsos positivos & $0.00658$ \\
Tasa de falsos negativos & $0.00465$ \\
\bottomrule
\end{tabular}
\caption[Métricas de la MAIN histórica]{Métricas de la ejecución MAIN histórica sobre su partición aleatoria fija. No son resultados de la campaña final. Fuente: \texttt{metrics.json} y \texttt{bootstrap\_ci\_seed42.json} históricos.}
\label{tab:main-metrics}
\end{table}

Las curvas de la Figura~\ref{fig:curvas-entrenamiento} documentan que la ejecución completó un proceso de aprendizaje coherente; no permiten inferir por sí mismas generalización ni estabilidad entre semillas.

\begin{figure}[htbp]
\centering
\includegraphics[width=0.92\textwidth]{f8_curvas_entrenamiento.pdf}
\caption[Curvas de entrenamiento de la MAIN histórica]{Curvas de entrenamiento de la ejecución MAIN histórica. Fuente: escalares TensorBoard exportados del artefacto previo a campaña.}
\label{fig:curvas-entrenamiento}
\end{figure}

\begin{figure}[htbp]
\centering
\includegraphics[width=0.64\textwidth]{f9_cm_main.pdf}
\caption[Matriz de confusión de la MAIN histórica]{Matriz de confusión de la MAIN histórica sobre la partición aleatoria fija. Fuente: \texttt{bootstrap\_ci\_seed42.json}; no es la matriz de la MAIN final.}
\label{fig:cm-main}
\end{figure}

\subsection{Precisión de muestreo y duplicados históricos}
\label{sec:resultados-bootstrap}

El \emph{bootstrap} histórico de $10\,000$ remuestreos estratificados sitúa la F1 de ataque en $[0.98394,0.98496]$, el \emph{recall} de ataque en $[0.99495,0.99575]$, la tasa de falsos positivos en $[0.00634,0.00681]$ y la tasa de falsos negativos en $[0.00425,0.00505]$. Estos intervalos cuantifican la precisión sobre un conjunto de prueba fijo y un único modelo; no estiman varianza de entrenamiento.

\begin{figure}[htbp]
\centering
\includegraphics[width=0.88\textwidth]{f10_ic_bootstrap.pdf}
\caption[Intervalos bootstrap históricos]{Intervalos al 95\% de la MAIN histórica sobre su test fijo. Fuente: \texttt{runs/validation/bootstrap\_ci\_seed42.json}. La campaña final regenerará este análisis.}
\label{fig:ic-bootstrap}
\end{figure}

El análisis histórico de duplicados encontró que el $22.30\%$ de las filas completas eran duplicados exactos y que el $24.86\%$ del test aleatorio coincidía con alguna fila de entrenamiento. Entre los ataques de prueba, la coincidencia ascendía al $40.12\%$. No demuestra por sí solo una fuga de etiquetas en la implementación, pero sí que la mezcla aleatoria ofrece una evaluación optimista sobre un benchmark con observaciones repetidas.

\begin{figure}[htbp]
\centering
\includegraphics[width=0.84\textwidth]{f12_duplicados.pdf}
\caption[Duplicados en la partición aleatoria histórica]{Duplicados exactos y coincidencias entre entrenamiento y prueba de la partición histórica con semilla~42. Fuente: \texttt{runs/validation/duplicate\_analysis\_seed42.json}.}
\label{fig:duplicados}
\end{figure}

\subsection{Controles A, B y proxy temporal histórico}
\label{sec:resultados-escalera}

La Comprobación~A reprodujo la predicción directa del modelo sobre $10\,000$ muestras y confirmó que las métricas no dependían del bucle del entorno. La Comprobación~B reentrenó brevemente con etiquetas barajadas y obtuvo $0.4773$ de exactitud frente a $0.5227$ de la clase mayoritaria, señal compatible con la ausencia de una vía trivial de fuga en ese control.

La Comprobación~C entrenó una red proxy \texttt{[512,256]} durante $30\,000$ pasos con lunes--miércoles y evaluó jueves--viernes. Registró $0.52954$ de \emph{recall} y $0.62578$ de F1 de ataque. Esta ejecución es una señal histórica de dificultad temporal, pero no sustituye el experimento final: usa otra arquitectura, otro presupuesto y un contrato de artefactos anterior.

\begin{figure}[htbp]
\centering
\includegraphics[width=0.64\textwidth]{f11_cm_check_c.pdf}
\caption[Matriz histórica del proxy por días]{Matriz de confusión de la Comprobación~C histórica. Fuente: \texttt{VAL\_checks\_C\_20260213\_004847/validation\_results.json}. Se reemplazará como evidencia temporal principal por el experimento completo de $3\,000\,000$ de pasos.}
\label{fig:cm-check-c}
\end{figure}

\subsection{Random Forest histórico como referencia interpretativa}
\label{sec:rf}

El barrido histórico de \emph{Random Forest} empleó la misma representación canónica y particiones equivalentes. En la partición aleatoria superó a QRDQN histórica en F1 de ataque ($0.99676$ frente a $0.98445$). En la partición por día, su \emph{recall} de ataque descendió a $0.08135$ y su F1 a $0.15005$, frente a $0.52954$ y $0.62578$ del proxy QRDQN. Esta comparación motiva una evaluación emparejada actual, pero no demuestra que QRDQN generalice mejor de forma global: el QRDQN temporal histórico no era MAIN y faltan escenarios repetidos.

\begin{figure}[htbp]
\centering
\includegraphics[width=0.92\textwidth]{f13_qrdqn_vs_rf.pdf}
\caption[Comparación histórica entre QRDQN y Random Forest]{Recall y F1 de ataque en las particiones histórica aleatoria y por día. La fila temporal de QRDQN corresponde a un proxy; no existe aquí una comparación final de campaña. Fuente: artefactos históricos MAIN, Check~C y RF.}
\label{fig:qrdqn-vs-rf}
\end{figure}

\begin{figure}[htbp]
\centering
\includegraphics[width=0.64\textwidth]{f14_cm_rf_dia.pdf}
\caption[Matriz histórica de Random Forest por día]{Matriz de confusión de Random Forest en la partición histórica por día. Fuente: \texttt{rf\_cicids2017\_canonical\_20260628\_024735\_\_day\_split/metrics.json}.}
\label{fig:cm-rf-dia}
\end{figure}

\subsection{Inferencia histórica de Fase 2}
\label{sec:fase2}

La ejecución histórica \path{P2v2_pred_20260610_161231_MAIN} aplicó la MAIN histórica a $2\,000\,000$ de flujos etiquetados del laboratorio del autor. Registró $0.991862$ de exactitud, $0.988452$ de \emph{recall} de ataque, $0.98380$ de F1 de ataque y una tasa de bloqueo de $0.252364$. La captura era cerrada, generada por el operador y no adversarial; por ello el resultado solo acredita que ese pipeline procesó esa captura concreta. La inferencia fresca de campaña volverá a ejecutar el análisis con la MAIN final y sus artefactos de preprocesamiento.

\begin{figure}[htbp]
\centering
\includegraphics[width=0.80\textwidth]{f15_fase2.pdf}
\caption[Resultados históricos de Fase 2]{Matriz y tasa de bloqueo de la inferencia histórica sobre tráfico de laboratorio. Fuente: \texttt{runs/phase2/P2v2\_pred\_20260610\_161231\_MAIN/metrics.json}. Su validez externa es limitada.}
\label{fig:fase2-resultados}
\end{figure}

\section{Inserción de la campaña final}
\label{sec:resultados-finales-pendientes}

La actualización final de este capítulo seguirá la Tabla~\ref{tab:estructura-resultados-finales}. La MAIN fresca se presentará junto con su validación directa, \emph{bootstrap}, duplicados y control de etiquetas; después se analizarán la partición completa por días, la curva de tamaño, la distribución entre semillas y los cuatro escenarios retenidos. Solo entonces se realizará la comparación emparejada con Random Forest y se cerrará con la inferencia fresca de Fase~2. Este orden separa capacidad intradistribución, robustez y transferencia, y evita repetir la misma métrica en varias secciones.

\evidencenote{No existen todavía valores para la MAIN fresca, el día completo, la escalera de tamaños, las semillas 42--46, los cuatro holdouts, los RF nuevos, la validación directa, el bootstrap, los duplicados, el control barajado ni la Fase~2 fresca. Las figuras cuantitativas finales se generarán únicamente cuando el agregador de campaña declare el conjunto completo.}

\section{Síntesis provisional}
\label{sec:resultados-sintesis}

La evidencia histórica permite afirmar que la implementación aprende la tarea definida y que el protocolo aleatorio, aun siendo útil como comprobación interna, no basta para sostener generalización. También justifica mantener Random Forest como línea base fuerte y tratar la transferencia al laboratorio como una comprobación offline de alcance limitado. No permite determinar todavía la estabilidad entre semillas, la eficiencia respecto al tamaño, la sensibilidad de los cuatro escenarios, el comportamiento del perfil completo por días ni el balance final entre QRDQN y RF. Esas son, exactamente, las preguntas reservadas a la campaña final.
